\section{Python SDK}


\subsection{Package Structure}

\begin{lstlisting}[caption=Python SDK Structure]
hanzo-python/
  hanzo/
    __init__.py
    client.py
    resources/
      __init__.py
      products.py
      orders.py
    models/
      __init__.py
      product.py
    errors.py
    utils/
      pagination.py
      webhooks.py
  setup.py
  pyproject.toml
\end{lstlisting}

\subsection{Client Implementation}

\begin{lstlisting}[language=Python, caption=Python Client]
import hanzo

client = hanzo.Client('sk_live_...')

# List products
products = client.products.list(limit=10)
for product in products:
    print(product.name)

# Create a product
product = client.products.create(
    name='Widget',
    price=1999,
    description='A useful widget',
)

# Async support
import asyncio

async def main():
    async_client = hanzo.AsyncClient('sk_live_...')
    products = await async_client.products.list()
    async for product in products:
        print(product.name)

asyncio.run(main())
\end{lstlisting}

\subsection{Type Annotations}

\begin{lstlisting}[language=Python, caption=Python Type Annotations]
from dataclasses import dataclass
from typing import Optional, List
from datetime import datetime

@dataclass
class Product:
    id: str
    name: str
    price: int
    description: Optional[str]
    variants: List['Variant']
    created_at: datetime
    updated_at: datetime

@dataclass
class ProductCreateParams:
    name: str
    price: int
    description: Optional[str] = None
\end{lstlisting}

\subsection{Context Managers}

\begin{lstlisting}[language=Python, caption=Python Context Manager]
with hanzo.Client('sk_live_...') as client:
    product = client.products.retrieve('prod_123')
    # Connection automatically closed
\end{lstlisting}

